\writeSectionFrame{\presentationTitle}{Das zu lösende Problem}

\begin{frame}{Unser Programm}
    \begin{center}
        \includegraphics[scale=0.45]{\BILDERPFAD/problem1.png}
    \end{center}
\end{frame}
\note{
    Wir möchten das folgende Programm. Konfigurationsdaten werden
    in einer properties-Datei abgelegt. Eine Configuration-Klasse
    liest diese Datei und schreibt die gelesenen Werte in eine Map.
    Dann kann die Konfiguration geändert werden und irgendwann werden
    die neuen Daten in die Datei zurückgeschrieben. In unserem Beispiel
    wird die Datei bei jedem Start neu angelegt, damit man die Effekte
    beobachten kann. 
}

\begin{frame}[fragile]{Ausgangsklasse}
	\begin{alertblock}{Configuration.java}
        \begin{lstlisting}
public final class Configuration {
	private static final String CONFIGURATION_FILE 
        = "configuration.properties";
	private Map<String, String> configurationMap;
	
    public Configuration() {
		initializeFile();
		configurationMap = new HashMap<String, String>();
		loadConfigurationFromFile();
	}       
        \end{lstlisting} 		
 	\end{alertblock}
\end{frame}

\begin{frame}[fragile]{Ausgangsklasse}
	\begin{alertblock}{Configuration.java}
        \begin{lstlisting}
private void loadConfigurationFromFile() {
    Properties properties = new Properties();
    File propertiesFile = new File(CONFIGURATION_FILE);
    
    try (FileInputStream fileInputStream = 
            new FileInputStream(propertiesFile)) {
        properties.load(fileInputStream);
        
        Set<String> keys = properties.stringPropertyNames();
        for (String key: keys) {
            String value = properties.getProperty(key);
            configurationMap.put(key, value);
        }
    } // catch-Zweige
}      
        \end{lstlisting} 		
 	\end{alertblock}
\end{frame}
\note{
    Hier wird die properties-Datei gelesen und die Werte werden in 
    die Map hineingeschrieben. 
}

\begin{frame}[fragile]{Ausgangsklasse}
	\begin{alertblock}{Configuration.java}
        \begin{lstlisting}
public void addValue(String key, String value) {
    if (configurationMap.containsKey(key)) {
        configurationMap.replace(key, value);
    } else {
        configurationMap.put(key, value);
    }
}

public Optional<String> getValue(String key) {
    String value = configurationMap.get(key);
    if (value == null) {
        return Optional.empty();
    } else {
        return Optional.of(value);
    }
}    
        \end{lstlisting} 		
 	\end{alertblock}
\end{frame}
\note{
    Die Methode addValue() schreibt einen neuen Wert in die Map, nicht
    aber in die Datei.
}

\begin{frame}[fragile]{Ausgangsklasse}
	\begin{alertblock}{Configuration.java}
        \begin{lstlisting}
public void writeToFile() {
    File file = new File(CONFIGURATION_FILE);

    try (FileWriter writer = new FileWriter(file);
         FileReader reader = new FileReader(file)) {

        Properties properties = new Properties();
        properties.load(reader);
        
        configurationMap.forEach((key, value) -> 
            properties.setProperty(key, value));

        LocalDateTime now = LocalDateTime.now();
        properties.store(writer, now.toString());
    } // catch-Zweige
}   
        \end{lstlisting} 		
 	\end{alertblock}
\end{frame}
\note{
    Hier wird die aktuelle Konfiguration in die Datei geschrieben. 
}

\begin{frame}[fragile]{Ausgangsklasse}
	\begin{alertblock}{ConfigurationMain.java}
        \begin{lstlisting}
private static void notSingletonExample() {
    Configuration configuration1 = new Configuration();
    configuration1.addValue("background", "black");
    configuration1.addValue("margin", "10mm");

    Configuration configuration2 = new Configuration();
    configuration2.addValue("italic", "true");
    configuration2.addValue("bold", "false");
    
    configuration1.writeToFile();
    configuration2.writeToFile();
}            
        \end{lstlisting}
 	\end{alertblock}
\end{frame}

\note{
 	Wir können jedoch mehrere Objekte erzeugen und haben somit mehrere Konfigurationen in unserem Programm,
 	die sich im Extremfall auch noch widersprechen. Unser Problem ist also, dass wir mehrere Objekte
 	von einer Klasse bilden können, obwohl dies unerwünscht ist. Behelfen kann man sich mit einer 
 	Programmierrichtlinie. Aber dann muss sich erstens jeder Programmierer dran halten und zweitens muss
 	man auch erkennen können, ob schon ein Objekt existiert. Besser wäre es, wenn man die Existenz nur
 	eines Objekts erzwingen könnte. Nach dem Ablauf dieses Programms werden wir nur eine einzige Konfiguration gespeichert; die andere ist verschwunden. 
}